% Complete Experimental Table - Docker vs Apple container Benchmark
% Automatically generated by analyze_results_enhanced.py
%
% To use in LaTeX, copy the contents of the desired section
%
\documentclass{article}
\usepackage{booktabs}
\usepackage{longtable}
\usepackage{multirow}
\usepackage{array}
\begin{document}

% ============================================
% TABLE: CPU/Memory
% ============================================

\begin{table}[h]
\centering
\small
\caption{Experimental Results: CPU/Memory}
\label{tab:cpu_memory}
\begin{tabular}{lccccccccc}
\toprule
Scenario & Tool & N & Mean & SD & Median & IQR & CI 95\% & Effect Size & p-value (FDR) \\
 & & & & & & & [lower, upper] & (Cohen's d) & \\
\midrule
\multirow{2}{*}{memory\_1t} & Docker & 30 & 5899.6533 & 2.3657 & 5899.6750 & 3.6425 & [5898.7700, 5900.5367] & -9.3289 & 0.0000*** \\
 & Container & 30 & 5931.0930 & 4.1375 & 5930.9950 & 2.6975 & [5929.5480, 5932.6380] &  &  \\
\midrule
\multirow{2}{*}{memory\_2t} & Docker & 30 & 9794.3543 & 2204.6995 & 10596.6950 & 358.3550 & [8971.1060, 10617.6026] & -0.5092 & 0.0051** \\
 & Container & 30 & 10589.0917 & 106.9912 & 10619.7800 & 22.1850 & [10549.1405, 10629.0428] &  &  \\
\midrule
\bottomrule
\end{tabular}
\begin{tablenotes}
\small
\item N: sample size; SD: standard deviation; IQR: interquartile range; CI: confidence interval; Effect Size: effect size (Cohen's d); p-value corrected by FDR (Benjamini-Hochberg).
\item Significance: *** p < 0.001, ** p < 0.01, * p < 0.05
\end{tablenotes}
\end{table}

\end{document}
